% Conclusion

\subsection{Summary}

This research evaluates CPU emulator accuracy for security-relevant x86-64 corner cases. During the second semester, the project expanded from preliminary probes to a broader comparison covering x87 status words, flags, timestamp behavior, arithmetic and SSE faults, MXCSR rounding, and machine-state instructions. The results show that every tested emulator differs from native hardware in at least one meaningful way.

\subsection{Implications}

Security researchers should validate emulator-based findings on native hardware whenever precise CPU behavior matters. In particular, fault delivery, status-word bits, and rounding-control behavior are strong correctness signals. Machine-state values are sometimes more environment-sensitive, but they remain valuable fingerprints for detecting emulated execution.

For emulator developers, the current probes provide compact regression tests. For security practitioners, the results provide practical guidance: Box64 currently has the strongest cluster of Linux user-mode correctness issues, QEMU TCG's main weakness is x87 status fidelity, Blink has clear flag and MXCSR mismatches, and shellcode-oriented backends vary substantially in stability and instruction support.

\subsection{Closing Thoughts}

The project now has two complementary tools: deterministic edge-case tests for known behaviors and \texttt{sandsifter\_rs} for broader instruction-space discovery. This combination is important because emulator accuracy problems are both known and unknown. Known problems require stable, explainable tests. Unknown problems require fuzzing, reduction, and careful interpretation.

As emulation becomes increasingly central to malware analysis, exploit development, and binary research, understanding its limitations becomes essential. The current work contributes a practical methodology for doing that: compare against native hardware, separate correctness bugs from fingerprints, and turn discovery findings into reproducible tests.
